% GENERATED FILE. Edit canonical lessons, then run npm run generate:books.
% canonical-source-hash: fnv1a64:eeaa68ac41d68997
% canonical-lessons: BN-C379-protishruti, BN-C379-songkolpo, BN-C379-drirho, BN-C379-bhonggo-kora, BN-C379-protigya-kora

\chapter{Keeping a Promise}
\label{ch:bn-keeping-a-promise}
\hlchaptermodality{379}

\begin{chapteropening}
\textbf{By the end of this chapter:} \emph{I can make a promise and say what happens if it is broken.}

Complete the last lesson of chapter 379: I can make a promise and say what happens if it is broken.
\end{chapteropening}

\section[\texorpdfstring{\bn{প্রতিশ্রুতি} (prôtishruti)}{প্রতিশ্রুতি (prôtishruti)}]{\bn{প্রতিশ্রুতি} (prôtishruti) — a promise; a pledge}

\label{lesson:BN-C379-protishruti}



\begin{quote}
Before the new one: say the Bengali for to postpone, then the Bengali for to chase up.
\end{quote}

\subsection*{You'll want to know: \bn{প্রতিশ্রুতি}}
\textbf{\bn{প্রতিশ্রুতি}} — \emph{prôtishruti} — \textquotedblleft{}a promise; a pledge\textquotedblright{}.

\begin{quote}
\bn{নেতা} \bn{যদি} \bn{প্রতিশ্রুতি} \bn{না} \bn{রাখেন}, \bn{তাহলে} \bn{আমরা} \bn{ভোট} \bn{দেব} \bn{না।}
\end{quote}

\begin{quote}
\emph{netā jôdi prôtishruti nā rākhen, tāhôle āmrā bhoṭ debô nā.}
\end{quote}

\begin{quote}
If the leader doesn't keep his promise, we won't vote for him.
\end{quote}

\begin{cousinweb}
\textbf{\bn{প্রতিশ্রুতি}} (\emph{prôtishruti}) is Sanskrit \emph{pratiśruti}, a promise (Monier-Williams), from \emph{prati-śru} 'to listen to, to agree'. It is given with \bn{দেওয়া} and kept with \bn{রাখা}.
\end{cousinweb}

\begin{grammarlens}[title={What you've built}]
One more word for saying what happens if.
\end{grammarlens}

\subsection*{Your turn}
\begin{itemize}
\raggedright
  \item \emph{Say it:} \emph{prôtishruti}
  \item \emph{Say it:} \emph{prôtishruti}, once more
  \item \emph{Say it:} \emph{netā jôdi prôtishruti nā rākhen, tāhôle āmrā bhoṭ debô nā.} — If the leader doesn't keep his promise, we won't vote for him.
  \item \emph{Recall:} say the Bengali for to postpone, then the Bengali for to chase up, then say \emph{prôtishruti} again
\end{itemize}

\begin{tcolorbox}[breakable,colback=teal!4,colframe=teal!35!black,title={Before you move on}]
What does \bn{প্রতিশ্রুতি} mean? (\textquotedblleft{}A promise; a pledge\textquotedblright{}.) Say it once more.
\end{tcolorbox}
\section[\texorpdfstring{\bn{সংকল্প} (sôngkôlpô)}{সংকল্প (sôngkôlpô)}]{\bn{সংকল্প} (sôngkôlpô) — a resolve; a firm intention}

\label{lesson:BN-C379-songkolpo}



\begin{quote}
Before the new one: say the Bengali for to chase up, then the Bengali for a pledge.
\end{quote}

\subsection*{You'll want to know: \bn{সংকল্প}}
\textbf{\bn{সংকল্প}} — \emph{sôngkôlpô} — \textquotedblleft{}a resolve; a firm intention\textquotedblright{}.

\begin{quote}
\bn{সংকল্প} \bn{থাকলে} \bn{কাজ} \bn{শেষ} \bn{হবেই।}
\end{quote}

\begin{quote}
\emph{sôngkôlpô thākle kāj shesh hôbei.}
\end{quote}

\begin{quote}
If you are determined, the work will get done.
\end{quote}

\begin{cousinweb}
\textbf{\bn{সংকল্প}} (\emph{sôngkôlpô}) is Sanskrit \emph{saṃkalpa}, a firm intention formed in the mind (Monier-Williams).
\end{cousinweb}

\begin{grammarlens}[title={What you've built}]
One more word for saying what happens if.
\end{grammarlens}

\subsection*{Your turn}
\begin{itemize}
\raggedright
  \item \emph{Say it:} \emph{sôngkôlpô}
  \item \emph{Say it:} \emph{sôngkôlpô}, once more
  \item \emph{Say it:} \emph{sôngkôlpô thākle kāj shesh hôbei.} — If you are determined, the work will get done.
  \item \emph{Recall:} say the Bengali for to chase up, then the Bengali for a pledge, then say \emph{sôngkôlpô} again
\end{itemize}

\begin{tcolorbox}[breakable,colback=teal!4,colframe=teal!35!black,title={Before you move on}]
What does \bn{সংকল্প} mean? (\textquotedblleft{}A resolve; a firm intention\textquotedblright{}.) Say it once more.
\end{tcolorbox}
\section[\texorpdfstring{\bn{দৃঢ়} (driṛhô)}{দৃঢ় (driṛhô)}]{\bn{দৃঢ়} (driṛhô) — firm; resolute}

\label{lesson:BN-C379-drirho}



\begin{quote}
Before the new one: say the Bengali for a pledge, then the Bengali for a resolve.
\end{quote}

\subsection*{You'll want to know: \bn{দৃঢ়}}
\textbf{\bn{দৃঢ়}} — \emph{driṛhô} — \textquotedblleft{}firm; resolute\textquotedblright{}.

\begin{quote}
\bn{তুমি} \bn{দৃঢ়} \bn{থাকলে} \bn{কেউ} \bn{তোমাকে} \bn{থামাতে} \bn{পারবে} \bn{না।}
\end{quote}

\begin{quote}
\emph{tumi driṛhô thākle keu tomāke thāmāte pārbe nā.}
\end{quote}

\begin{quote}
If you stand firm, no one can stop you.
\end{quote}

\begin{cousinweb}
\textbf{\bn{দৃঢ়}} (\emph{driṛhô}) is Sanskrit \emph{dṛḍha}, fixed and firm (Monier-Williams); the ḍh between vowels is written \bn{ঢ়} and said ṛh.
\end{cousinweb}

\begin{grammarlens}[title={What you've built}]
One more word for saying what happens if.
\end{grammarlens}

\subsection*{Your turn}
\begin{itemize}
\raggedright
  \item \emph{Say it:} \emph{driṛhô}
  \item \emph{Say it:} \emph{driṛhô}, once more
  \item \emph{Say it:} \emph{tumi driṛhô thākle keu tomāke thāmāte pārbe nā.} — If you stand firm, no one can stop you.
  \item \emph{Recall:} say the Bengali for a pledge, then the Bengali for a resolve, then say \emph{driṛhô} again
\end{itemize}

\begin{tcolorbox}[breakable,colback=teal!4,colframe=teal!35!black,title={Before you move on}]
What does \bn{দৃঢ়} mean? (\textquotedblleft{}Firm; resolute\textquotedblright{}.) Say it once more.
\end{tcolorbox}
\section[\texorpdfstring{\bn{ভঙ্গ} \bn{করা} (bhônggô kôrā)}{ভঙ্গ করা (bhônggô kôrā)}]{\bn{ভঙ্গ} \bn{করা} (bhônggô kôrā) — to break a promise; to break (a rule)}

\label{lesson:BN-C379-bhonggo-kora}



\begin{quote}
Before the new one: say the Bengali for a resolve, then the Bengali for resolute.
\end{quote}

\subsection*{You'll want to know: \bn{ভঙ্গ} \bn{করা}}
\textbf{\bn{ভঙ্গ} \bn{করা}} — \emph{bhônggô kôrā} — \textquotedblleft{}to break a promise; to break (a rule)\textquotedblright{}.

\begin{quote}
\bn{প্রতিশ্রুতি} \bn{ভঙ্গ} \bn{করলে} \bn{মানুষ} \bn{আর} \bn{বিশ্বাস} \bn{করে} \bn{না।}
\end{quote}

\begin{quote}
\emph{prôtishruti bhônggô kôrle mānush ār bishvās kôre nā.}
\end{quote}

\begin{quote}
If you break a promise, people stop trusting you.
\end{quote}

Past (he/she …): \bn{তিনি} \bn{চুক্তি} \bn{ভঙ্গ} \bn{করলেন}, \emph{tini chukti bhônggô kôrlen}.

\begin{cousinweb}
\textbf{\bn{ভঙ্গ} \bn{করা}} (\emph{bhônggô kôrā}) is Sanskrit \emph{bhaṅga}, breaking (Monier-Williams), with \bn{করা}. It is the formal word for breaking a promise, a law or a silence.
\end{cousinweb}

\begin{grammarlens}[title={What you've built}]
One more word for saying what happens if.
\end{grammarlens}

\subsection*{Your turn}
\begin{itemize}
\raggedright
  \item \emph{Say it:} \emph{bhônggô kôrā}
  \item \emph{Say it:} \emph{bhônggô kôrā}, once more
  \item \emph{Say it:} \emph{prôtishruti bhônggô kôrle mānush ār bishvās kôre nā.} — If you break a promise, people stop trusting you.
  \item \emph{Recall:} say the Bengali for a resolve, then the Bengali for resolute, then say \emph{bhônggô kôrā} again
\end{itemize}

\begin{tcolorbox}[breakable,colback=teal!4,colframe=teal!35!black,title={Before you move on}]
What does \bn{ভঙ্গ} \bn{করা} mean? (\textquotedblleft{}To break a promise; to break (a rule)\textquotedblright{}.) Say it once more.
\end{tcolorbox}
\section[\texorpdfstring{\bn{প্রতিজ্ঞা} \bn{করা} (prôtigyā kôrā)}{প্রতিজ্ঞা করা (prôtigyā kôrā)}]{\bn{প্রতিজ্ঞা} \bn{করা} (prôtigyā kôrā) — to vow; to swear (to do something)}

\label{lesson:BN-C379-protigya-kora}



\begin{quote}
Before the new one: say the Bengali for resolute, then the Bengali for to break a promise.
\end{quote}

\subsection*{You'll want to know: \bn{প্রতিজ্ঞা} \bn{করা}}
\textbf{\bn{প্রতিজ্ঞা} \bn{করা}} — \emph{prôtigyā kôrā} — \textquotedblleft{}to vow; to swear (to do something)\textquotedblright{}.

\begin{quote}
\bn{আমি} \bn{প্রতিজ্ঞা} \bn{করলাম} \bn{যে} \bn{কাজ} \bn{শেষ} \bn{না} \bn{হলে} \bn{আমি} \bn{ছুটি} \bn{নেব} \bn{না।}
\end{quote}

\begin{quote}
\emph{āmi prôtigyā kôrlām je kāj shesh nā hôle āmi chhuṭi nebô nā.}
\end{quote}

\begin{quote}
I vowed that I would not take a holiday until the work was done.
\end{quote}

Past (he/she …): \bn{তিনি} \bn{সত্যি} \bn{বলার} \bn{প্রতিজ্ঞা} \bn{করলেন}, \emph{tini sôtti bôlār prôtigyā kôrlen}.

\begin{cousinweb}
\textbf{\bn{প্রতিজ্ঞা} \bn{করা}} (\emph{prôtigyā kôrā}) is Sanskrit \emph{pratijñā}, a vow (Monier-Williams), with \bn{করা}; \bn{জ্ঞ} is said gy. It is a solemn promise to yourself or others.
\end{cousinweb}

\begin{grammarlens}[title={What you've built}]
One more word for saying what happens if.
\end{grammarlens}

\subsection*{Your turn}
\begin{itemize}
\raggedright
  \item \emph{Say it:} \emph{prôtigyā kôrā}
  \item \emph{Say it:} \emph{prôtigyā kôrā}, once more
  \item \emph{Say it:} \emph{āmi prôtigyā kôrlām je kāj shesh nā hôle āmi chhuṭi nebô nā.} — I vowed that I would not take a holiday until the work was done.
  \item \emph{Recall:} say the Bengali for resolute, then the Bengali for to break a promise, then say \emph{prôtigyā kôrā} again
\end{itemize}

\begin{tcolorbox}[breakable,colback=teal!4,colframe=teal!35!black,title={Before you move on}]
What does \bn{প্রতিজ্ঞা} \bn{করা} mean? (\textquotedblleft{}To vow; to swear (to do something)\textquotedblright{}.) Say it once more.
\end{tcolorbox}
